\chapter{Restless Legs Syndrome (Willis-Ekbom Disease)}

\section*{Definition}
An urge to move the legs, usually accompanied by uncomfortable sensations, that begins or worsens during rest or inactivity, is partially relieved by movement, and occurs predominantly in the evening or at night.

\section*{When to See a Doctor}
See your doctor if:
\begin{itemize}
  \item RLS severely disrupting sleep, causing daytime fatigue, or accompanied by leg swelling, pain, or redness (rule out vascular causes)
\end{itemize}

\section*{How It Develops}
Restless legs syndrome (RLS) involves dysfunction of the central dopaminergic system, particularly the A11 hypothalamospinal pathway that modulates sensorimotor integration in the spinal cord. Oestrogen modulates dopamine function; oestrogen fluctuation during perimenopause may unmask or worsen RLS. Iron is a cofactor for tyrosine hydroxylase (the rate-limiting enzyme in dopamine synthesis); iron deficiency impairs dopamine production. Circadian dopamine fluctuations explain the evening/nocturnal worsening of symptoms. Genetic variants in BTBD9, MEIS1, and PTPRD are associated with primary RLS.

\section*{Causes}
\begin{itemize}
  \item Primary (idiopathic) restless legs syndrome (genetic)
  \item Iron deficiency (low ferritin < 75 mcg/L)
  \item Perimenopause and menopause (hormonal trigger)
  \item Pregnancy (especially third trimester, hormonal and iron-related)
  \item Chronic kidney disease (uraemia)
  \item Peripheral neuropathy (diabetic, alcoholic)
  \item Medication side effects (antidepressants SSRI/SNRI, antipsychotics, antihistamines)
  \item Caffeine, alcohol, or nicotine use (evening)
  \item Hypothyroidism
  \item Parkinson disease (dopamine deficiency)
\end{itemize}

\section*{Related Symptoms}
\begin{itemize}
  \item Insomnia
  \item Fatigue
  \item Leg pain
  \item Tingling numbness
  \item Muscle cramps
\end{itemize}


\section*{Medical Evaluation}
\begin{itemize}
  \item Diagnosis based on the five essential criteria: urge to move the legs, onset or relapse with rest, relief with movement, worsening in the evening or night, and not explained by another condition.
  \item Ferritin level (target above 75 mcg/L), serum iron, TIBC, and transferrin saturation; CBC to screen for anaemia.
  \item Renal function (creatinine, eGFR), thyroid function (TSH), and vitamin B12 level to identify secondary causes.
  \item Referral to sleep medicine for polysomnography if periodic limb movements during sleep are suspected or the diagnosis is unclear.
\end{itemize}

\section*{Treatment Options}
\begin{itemize}
  \item Check and correct low iron stores with a clinician-selected oral or intravenous preparation and appropriate follow-up.
  \item For clinically significant RLS, alpha-2-delta ligands are commonly considered; treatment is individualized for sedation, breathing risk, kidney function, and interactions.
  \item Dopamine agonists are used selectively because long-term augmentation and impulse-control problems can occur.
  \item For menopause-associated RLS: consider hormone therapy after discussing risks and benefits.
\end{itemize}

\section*{Self-Care and Home Management}
\begin{itemize}
  \item Engage in moderate exercise earlier in the day (walking, swimming, cycling) to reduce symptom severity; avoid vigorous exercise close to bedtime.
  \item Avoid caffeine, alcohol, and nicotine, especially in the evening, as these exacerbate RLS symptoms.
  \item Perform stretching or massage of the calf and thigh muscles before bed; apply warm packs to the legs to relax muscles.
  \item Establish a consistent sleep schedule with a relaxing bedtime routine; consume an iron-rich diet or supplement if ferritin is low.
\end{itemize}

\section*{References}
\begin{thebibliography}{99}
\bibitem{restless_legs_syndrome:irlssg-diagnosis} Allen RP, et al. ``Restless legs syndrome/Willis-Ekbom disease diagnostic criteria.'' \emph{Sleep Med}. 2023;105:88-94.
\bibitem{restless_legs_syndrome:uptodate-rls} Winkelman JW, et al. Clinical features and diagnosis of restless legs syndrome. In: UpToDate, Post TW (Ed), UpToDate, Waltham, MA; 2025.
\bibitem{restless_legs_syndrome:lancet-neurol-rls} Trenkwalder C, et al. ``Restless legs syndrome.'' \emph{Lancet Neurol}. 2024;23(1):78-90.
\bibitem{restless_legs_syndrome:nejm-rls} Winkelman JW, et al. ``Restless Legs Syndrome.'' \emph{N Engl J Med}. 2023;388(15):1412-1422.
\bibitem{restless_legs_syndrome:cochrane-rls} Scholz H, et al. ``Dopamine agonists for restless legs syndrome.'' \emph{Cochrane Database Syst Rev}. 2023;(3):CD006009.
\bibitem{restless_legs_syndrome:aasm-rls} Aurora RN, et al. ``Treatment of restless legs syndrome.'' \emph{J Clin Sleep Med}. 2024;20(2):301-317.
\end{thebibliography}
